\chapter{Normal Domains of Attraction}
\label{sec:p15-normal-domains-of-attraction}

Finite variance is sufficient for a normal limit, but it is not
necessary when the normalization may grow faster than \(\sqrt n\).
The exact criterion compares the probability of a large observation
with the accumulated second moment below that observation's scale.
This chapter proves both directions of the criterion stated in
Durrett's Theorem 3.4.14
\cite[p.~152]{DurrettPTEAuthor2019}.
The first proof is an independent construction from the supplied
probability foundations. A separate section retains the cutoff
argument of the acquired 1954 Gnedenko--Kolmogorov presentation
\cite[Section~35, Theorem~1, pp.~172--175]{GnedenkoKolmogorov1954LimitDistributions}.

\begin{hbdefinition}{def:p15-nda-law-and-moments}{Attraction, tails, and truncated moments}
Let \(X_1,X_2,\ldots\) be i.i.d. finite real random variables,
with common law that is nondegenerate, meaning not concentrated
at any single point. Write \(X=X_1\) and
\(S_n=\sum_{j=1}^nX_j\).
The law belongs to the \emph{domain of attraction of the normal law}
if there are deterministic real \(a_n\) and positive finite \(b_n\)
such that
\[
 (S_n-a_n)/b_n\Rightarrow\chi,
 \]
where \(\chi\) has density
\((2\pi)^{-1/2}e^{-x^2/2}\).
Write \(N(0,\tau^2)\) for the law of \(\tau\chi\), with
\(\tau>0\). A symmetric law is one unchanged by replacing
its variable by its negative.
The word \emph{normal} describes the limiting distribution; it does
not stipulate that the scale is a fixed multiple of \(\sqrt n\).

For any finite real random variable \(Z\), set, for \(x>0\),
\[
 q_Z(x)=\mathbb P(|Z|>x),\qquad
 H_Z(x)=\mathbb E(Z^2;|Z|\leq x),\qquad
 W_Z(x)=\mathbb E\min(Z^2,x^2)=H_Z(x)+x^2q_Z(x).
 \]
At zero define \(q_Z(0)=\mathbb P(|Z|>0)\), when a tail
quantile uses that endpoint.
These truncated moments are finite. They become strictly positive
for all sufficiently large \(x\) if \(\mathbb P(Z\ne0)>0\).
A positive function \(H\), defined for all sufficiently large
positive arguments, is \emph{slowly varying} at infinity if
\(H(cx)/H(x)\to1\) for every fixed \(c>0\).
Whenever a ratio has \(H_Z(x)\) in its denominator it is used
only on this eventually positive domain.
All logarithms are natural; \(f(x)\sim g(x)\) means
\(f(x)/g(x)\to1\), and \(f=o(g)\) means \(f/g\to0\).
\end{hbdefinition}

\section{Tail and moment equivalences}
\label{sec:p15-nda-tail-calculus}

\begin{hbtheorem}{thm:p15-nda-tail-equivalence}{proved=annex:p15.nda-tail-equivalence}{Slow variation is exactly the negligible-tail condition}
Let \(Z\) be finite and \(\mathbb P(Z\ne0)>0\). Then
\[
 \frac{x^2q_Z(x)}{H_Z(x)}\longrightarrow0
       \quad\Longleftrightarrow\quad H_Z\text{ is slowly varying}.
 \]
When these conditions hold, \(W_Z(x)\sim H_Z(x)\) and is
slowly varying as well. Also
\[
 \frac{H_Z(x-)}{H_Z(x)}\to1,\qquad
 \frac{x^2\mathbb P(|Z|\geq x)}{H_Z(x)}\to0.
 \]
Here \(H_Z(x-)=\mathbb E(Z^2;|Z|<x)\).
Thus strict or closed truncation conventions yield the same
asymptotic criterion, including laws with atoms.
\end{hbtheorem}

\begin{hbtheorem}{thm:p15-nda-moments-and-shifts}{proved=annex:p15.nda-moments-and-shifts}{Moment bounds, integrable tails, and centering}
Under the equivalent conditions above, \(\mathbb E|Z|^p<\infty\)
for every \(0<p<2\), and
\[
 \mathbb E(|Z|^p;|Z|>x)
                    =o\{H_Z(x)/x^{2-p}\}.
 \]
In particular the mean exists absolutely.
If the law is nondegenerate and \(c\in\mathbb R\) is fixed,
the same tail criterion and slow-variation condition hold for \(Z-c\).
If \(\mathbb EZ^2=\infty\), then
\(H_{Z-c}(x)\sim H_Z(x)\).
If \(\mathbb EZ^2<\infty\), both truncated moments tend to
their respective positive full second moments, which need not be equal.
For \(Y=X-\mu\), where \(\mu=\mathbb EX\),
write \(V(x)=H_Y(x)\); its finite-moment limit is
\(\operatorname{Var}(X)\), not \(\mathbb EX^2\).
\end{hbtheorem}

\section{Construction and sufficiency}
\label{sec:p15-nda-sufficiency}

\begin{hbtheorem}{thm:p15-nda-normalizer-construction}{proved=annex:p15.nda-normalizer-construction}{A continuous construction of every required scale}
Suppose \(Y\) is centered, nondegenerate, and \(V=H_Y\) is
slowly varying. Choose \(x_0>0\) with
\(\mathbb P(0<|Y|\leq x_0)>0\).
For all sufficiently large integers \(n\), there is a unique
\(b_n\geq x_0\) satisfying
\[
 \frac{W_Y(b_n)}{b_n^2}=\frac1n.
 \]
Choose any positive finite \(b_n\) at the finitely many earlier
indices. Then \(b_n\to\infty\), it is increasing eventually, and
\[
 \frac{nV(b_n)}{b_n^2}\to1,\qquad
 n q_Y(\varepsilon b_n)\to0,\qquad
 \frac{n}{b_n^2}V(\varepsilon b_n)\to1
                                  \quad(\varepsilon>0).
 \]
Moreover
\[
 \frac{n}{b_n}\mathbb E(|Y|;|Y|>b_n)\to0,\qquad
 \frac{b_{\lfloor cn\rfloor}}{b_n}\to\sqrt c\quad(c>0).
 \]
The equation uses the continuous winsorized moment \(W_Y\);
it does not assert that a discontinuous truncated-moment equation
has an exact solution.
\end{hbtheorem}

\begin{hbtheorem}{thm:p15-nda-sufficiency}{proved=annex:p15.nda-sufficiency}{The full infinite- or finite-variance sufficient condition}
If the nondegenerate i.i.d. law of \(X\) satisfies
\[
 x^2q_X(x)/H_X(x)\to0,
 \]
let \(\mu=\mathbb EX\), \(Y=X-\mu\), and choose \(b_n\)
by the preceding continuous construction. Then
\[
 \frac{S_n-n\mu}{b_n}\Rightarrow\chi.
 \]
The proof verifies the Lindeberg condition for the full centered
truncated array, proves the vanishing probability of any removed
term, and proves that truncated-mean centering may be replaced by
the exact center \(n\mu\).
\end{hbtheorem}

\section{Necessity without an assumed infinite-divisibility theorem}
\label{sec:p15-nda-necessity}

\begin{hblemma}{thm:p15-nda-normal-types}{proved=annex:p15.nda-normal-types}{Uniqueness of affine normalizations}
If \(Z_n\Rightarrow\chi\), \(c_n>0\), and
\(c_nZ_n+d_n\Rightarrow\chi\) for deterministic real \(d_n\),
then \(c_n\to1\) and \(d_n\to0\).
The same assertion holds with the same nondegenerate centered
normal law in both limits.
\end{hblemma}

\begin{hbtheorem}{thm:p15-nda-symmetric-necessity}{proved=annex:p15.nda-symmetric-necessity}{A symmetric normal limit forces negligible jumps at every scale}
Let \(Z_1,Z_2,\ldots\) be i.i.d., symmetric, and nondegenerate.
Suppose \(b_n>0\) and
\[
 \frac{Z_1+\cdots+Z_n}{b_n}\Rightarrow N(0,\tau^2),
 \qquad 0<\tau^2<\infty.
 \]
Then \(b_n\to\infty\), \(b_{n+1}/b_n\to1\), and for every \(r>0\),
\[
 n q_Z(rb_n)\to0,\qquad
 \frac{n}{b_n^2}H_Z(rb_n)\to\tau^2.
 \]
Consequently \(H_Z\) is slowly varying on the whole positive
half-line at infinity, and \(x^2q_Z(x)/H_Z(x)\to0\).
The proof supplies the nonnegative trigonometric-kernel argument
and the passage from the sequence \(b_n\) to arbitrary \(x\).
The scales \(b_n\) need not initially be monotone.
\end{hbtheorem}

\begin{hbtheorem}{thm:p15-nda-desymmetrization}{proved=annex:p15.nda-desymmetrization}{Transferring the criterion from an independent difference}
Let \(X'\) be an independent copy of a nondegenerate finite
real \(X\), and put \(Z=X-X'\).
If \(H_Z\) is slowly varying, then
\[
 x^2q_X(x)/H_X(x)\to0.
 \]
No expectation or second moment of \(X\) is assumed in this
implication; the required integrability follows from its conclusion.
\end{hbtheorem}

\begin{hbtheorem}{thm:p15-nda-characterization}{proved=annex:p15.nda-characterization}{The complete normal-domain characterization}
For a nondegenerate i.i.d. law on \(\mathbb R\), the following
conditions are equivalent:
\begin{enumerate}
\item There are deterministic \(a_n\in\mathbb R\), \(d_n>0\)
 with \((S_n-a_n)/d_n\Rightarrow\chi\).
\item \(x^2q_X(x)/H_X(x)\to0\).
\item \(H_X\) is slowly varying.
\end{enumerate}
In this case \(\mu=\mathbb EX\) exists absolutely. Let \(b_n\)
be the constructed scales for \(X-\mu\). For any deterministic
\(a_n\in\mathbb R\), \(d_n>0\),
\[
 \frac{S_n-a_n}{d_n}\Rightarrow\chi
 \quad\Longleftrightarrow\quad
 \frac{d_n}{b_n}\to1
 \quad\text{and}\quad
 \frac{a_n-n\mu}{b_n}\to0.
 \]
Thus the construction also describes every normalization giving
this prescribed standard normal limit.
\end{hbtheorem}

\section{Examples and sharp boundaries}
\label{sec:p15-nda-examples}

\begin{hbtheorem}{thm:p15-nda-examples}{proved=annex:p15.nda-examples}{Finite variance, logarithmic variance, and excluded power tails}
\begin{enumerate}
\item If \(0<\operatorname{Var}(X)=\sigma^2<\infty\), then
 the constructed \(b_n\sim\sigma\sqrt n\).
 Using \(\sqrt{n\mathbb EX^2}\) instead, when the mean is
 nonzero, yields limiting variance
 \(\sigma^2/\mathbb EX^2\), not one.
\item For the symmetric law with density
 \(|x|^{-3}\mathbf1_{\{|x|>1\}}\),
 \[
 H_X(x)=2\log x,\quad W_X(x)=1+2\log x\quad(x\geq1),
 \qquad b_n\sim\sqrt{n\log n}.
 \]
 Its variance is infinite, yet \(S_n/\sqrt{n\log n}\Rightarrow\chi\).
\item For a symmetric law with
 \(\mathbb P(|X|>x)=x^{-\alpha}\), \(x\geq1\), where
 \(0<\alpha<2\),
 \[
 H_X(x)=\frac{\alpha}{2-\alpha}(x^{2-\alpha}-1),\qquad
 \frac{x^2q_X(x)}{H_X(x)}\to\frac{2-\alpha}{\alpha}>0.
 \]
 No deterministic centering and positive normalization can make
 its i.i.d. sums converge to a nondegenerate normal law.
\item Nondegeneracy cannot be removed. If \(X=1\) almost surely,
 the tail ratio is zero for \(x>1\), while every normalized sum
 is deterministic and cannot have a nondegenerate normal limit.
 If \(X=0\) almost surely, the displayed ratio has zero denominator
 and is not a defined condition.
\end{enumerate}
\end{hbtheorem}

\section{The separate 1954 cutoff presentation}
\label{sec:p15-nda-gk-source}

\begin{hblemma}{thm:p15-nda-gk-truncated-mean}{proved=annex:p15.nda-gk-truncated-mean}{The source's weighted Cauchy--Schwarz premise}
For any finite real \(X\) with \(\mathbb EX^2=\infty\), set
\(m(C)=\mathbb E(X;|X|\leq C)\), \(U(C)=H_X(C)\).
Then, as \(C\to\infty\),
\[
 m(C)^2=o(U(C)),\qquad
 \operatorname{Var}(X\mathbf1_{\{|X|\leq C\}})\sim U(C).
 \]
No first moment of \(X\) is assumed here.
The proof constructs the positive weight used on source p.~173
and proves its square integrability and growth at both infinities.
\end{hblemma}

\begin{hbtheorem}{thm:p15-nda-gk-cutoff-interface}{proved=annex:p15.nda-gk-cutoff-interface}{The exact i.i.d. cutoff interface used by the source}
For a nondegenerate i.i.d. law with \(\mathbb EX^2=\infty\),
normal attraction is equivalent to the existence of deterministic
\(C_n\to\infty\) with
\[
 n q_X(C_n)\to0,\qquad
 \frac n{C_n^2}
 \operatorname{Var}(X\mathbf1_{\{|X|\leq C_n\}})\to\infty.
 \]
Equivalently, the second condition is \(nU(C_n)/C_n^2\to\infty\).
Given these cutoffs, put
\[
 B_n^2=n\operatorname{Var}(X\mathbf1_{\{|X|\leq C_n\}}),
 \qquad a_n=n m(C_n).
 \]
These quantities have \(B_n>0\) eventually and
\((S_n-a_n)/B_n\Rightarrow\chi\).
This is the infinite-variance i.i.d. specialization of the
source's Section 26, Theorem 4, which its Section 35 proof actually
uses. The annex proves this whole specialization; it does not
claim the broader non-identically-distributed theorem.
\end{hbtheorem}

\begin{hbtheorem}{thm:p15-nda-gk-cutoff-route}{proved=annex:p15.nda-gk-cutoff-route}{The source's cutoff comparison and tail-quantile construction}
Write \(U(x)=H_X(x)\), \(\mathcal H(x)=U(x)/x^2\), and
\(q(x)=q_X(x)\). If \(\mathbb EX^2=\infty\), then
\[
 \frac{q(x)}{\mathcal H(x)}\to0
 \quad\Longleftrightarrow\quad
 \text{there are }C_n\to\infty
 \text{ with }nq(C_n)\to0,\quad n\mathcal H(C_n)\to\infty.
 \]
For the constructive direction, for fixed \(\delta>0\) set
\[
 C_n(\delta)=\inf\{x\geq0:nq(x)\leq\delta\}.
 \]
For each such \(\delta\), these cutoffs tend to infinity and
\(n\mathcal H(C_n(\delta))\to\infty\).
A fully specified diagonal choice \(\delta_n\downarrow0\) gives
the desired \(C_n=C_n(\delta_n)\).
The preceding cutoff interface then proves normal attraction
by the source's distinct route. The finite-variance case is
supplied separately in the proof, as in source pp.~172--173.
The retained event is closed and the removed tail is strict;
the already proved endpoint equivalence supplies the source's
open-truncation version of the final ratio criterion.
\end{hbtheorem}

\begin{proofboundary}{def:p15-nda-historical-boundary}{Completed normal criterion and the broader historical chain}
Both directions of Durrett's Theorem 3.4.14 and the separate
Section 35, Theorem 1 cutoff argument in the 1954 source are
supplied, with their complete i.i.d. interface.
The independent trigonometric-kernel proof of necessity is
distinguished from that source argument.
The broader general-array infinitely divisible-law presentations
in Gnedenko--Kolmogorov Sections 19--26, including their own
canonical-measure proofs, remain separate presentations. The
i.i.d. specialization proved here does not count those larger
historical theorem scopes as reproduced.
\end{proofboundary}
